\clearpage
\fancyhead[L]{\small CONNAISSANCE DES TEMS.}
\SourceLeaf{532}
\begin{center}\small Göttingische gelehrte Anzeigen. Stück 153. Seite 1526..1528. 1810 September 24.\\[3pt]\rule{.72\linewidth}{.3pt}\end{center}

\emph{Connaissance des tems ou des mouvemens célestes, à l'usage des astronomes et des navigateurs, pour l'an} 1811, \emph{publiée par le bureau des longitudes.} 1809. Octav. Paris.

Wir übergehen den Kalender, dessen Einrichtung unverändert geblieben ist, um nur den Inhalt der Zusätze anzuzeigen. Drei Viertel davon machen die astronomischen, auf der kaiserl. Sternwarte zu Paris in den Jahren 1807 und 1808 angestellten, Beobachtungen aus, deren Wichtigkeit wir hier nicht erst zu rühmen brauchen. Ausser der Sonne, dem Monde und den vornehmsten Fixsternen wurden die sämmtlichen ältern Planeten sehr fleissig beobachtet; von den neuen finden sich im Jahre 1807 auch zahlreiche Beobachtungen, besonders von der Vesta und Ceres: die Pallas wurde verschiedene Male beobachtet, Juno gar nicht. Im Jahre 1808 finden wir von den neuen Planeten blos Vesta und Juno beobachtet, letztere nur Ein Mal. Auffallend sind uns in diesem Beobachtungs-Journale die zuweilen Statt findenden sehr beträchtlichen Ungleichheiten im Gange der Pendeluhr, denen man durch öfteres Verrücken der Linse abzuhelfen suchte. An der parallatischen Maschine wurden die Vesta und der grosse Comet von 1807 beobachtet; von letzterem theilt \textsc{Bouvard} seine Elemente mit. Hierauf folgt ein Aufsatz von \textsc{Laplace} über die Abnahme der Schiefe der Ekliptik, nach alten Beobachtungen. Besonders interessant ist hier die Erörterung einiger alten Chinesischen Beobachtungen, wovon die älteste 1100 Jahre vor unserer Zeitrechnung gemacht ist, und welche alle eben so, wie die Griechischen, Arabischen und Persischen Bestimmungen, mit der von \textsc{Laplace} aus der Theorie bestimmten Abnahme so gut zusammentreffen, als sich nur immer erwarten lässt. Ein anderer Aufsatz von \textsc{Laplace} betrifft die Rotation des Saturnringes, worüber bekanntlich \textsc{Schröter}'s Beobachtungen ein von den \textsc{Herschel}'schen ganz verschiedenes und mit der Theorie unvereinbares Resultat gegeben haben. \textsc{Laplace} erklärt die erstern, nach welchen auf dem Ringe Ungleichheiten zu sein schienen, die ihren Platz nicht änderten, durch die sinnreiche Hypothese, dass der Ring eigentlich aus vielen concentrischen, aber in verschiedenen Ebenen liegenden, Ringen bestehe, die allerdings solche Ungleichheiten zeigen konnten. So lange die Ebenen der einzelnen Ringe ihre Lage nicht merklich änderten, würden diese Ungleichheiten freilich, wenn sie gleich rotirten, unverändert erscheinen; nur bleibt hierbei noch unerklärt, warum beide Ansen nicht einerlei Anblick darboten. Man darf hoffen, dass die künftige Wiederholung dieser delicaten Beobachtungen hierüber Aufschluss geben wird. — Eine hierauf folgende Abhandlung von \textsc{Delambre} über die Breiten- und Zeitbestimmung aus den beobachteten Höhen zweier bekannter Sterne, ist hauptsächlich zu einem Commentar über ein auch in unsern Blättern angezeigtes Programm von \textsc{Gauss} (s. Gött. gel. Anz. 1808 S. 1945 [S. 37 d. B.]) bestimmt, obwohl \textsc{Delambre} diese kleine Schrift nur aus dem in der Monatlichen Correspondenz befindlichen Auszuge gekannt zu haben scheint, und daher den Zweck derselben nicht ganz richtig beurtheilt. Auch die in der Monatl. Correspondenz von \textsc{Mollweide} mitgetheilte Auflösung desselben Problems führt \textsc{Delambre} an und rühmt an ihr, obwohl er sie aus andern Rücksichten der trigonometrischen nachsetzt, die Leichtigkeit, womit bei derselben von den beiden Wurzeln der Gleichung die rechte ausgewählt werden könne. Diesem Urtheil können wir nicht beipflichten: denn das von \textsc{Mollweide} angeführte Kriterium ist zwar einfach, aber unstatthaft, und der Aufgabe fremd. — Berechnung der Beobachtungen der Sonnenfinsterniss vom 16. Juni 1806 zu Utrecht, Mailand, München und Lilienthal von \textsc{van Beek Calkoen}. — \textsc{Burckhardt} über die verschiedenen, von den Astronomen in Anwendung gebrachten, Arten, die Sonne zu beobachten. — Von demselben Astronomen noch drei kleine Aufsätze über die Cometen von 1701 und 1772, und über das rostförmige Pen\SourceLeaf{533}del. — Von Werken, welche in die Astronomie einschlagen, ist diesmal blos die Reise von \textsc{d'Entrecasteaux} angezeigt. — Ein Auszug aus dem auf der Sternwarte von \textsc{Bouvard} im Jahre 1807 geführten meteorologischen Journal macht den Beschluss.
